% Replication programme (November 2026). Numbers: docs/REPLICATION_RESULTS.md (final) and results/replication/.
\section{Validation against published methods (replication programme)}\label{sec:replication}
\plain{To make sure our conclusions are not just about methods we built ourselves, we took eight published papers on
shortcut learning in medical images, ran their best method with the authors' own code, and compared it with ours ---
first on \emph{their} data and test, then on \emph{our} real-artifact test (thyroid calipers). Every rule was written
down before the results. The main outcome: published methods that restrict the model to the region of interest fail
when the artifact is inside it, and our class-conditional constraint beats ISNet, RR-ClArC and EPVT under registered
rules.}

\subsection{How the comparison works}
\begin{itemize}[leftmargin=*]
\item \textbf{Reproduction gate.} We first ran each paper's best method with its code and compared it with the published
number (tolerance 0.05 in the paper's metric; widened from 0.03 for every paper after the first gates were known --- a
deviation we report).
\item \textbf{Tier 1 (reproduced):} a formal ``beaten / not beaten'' verdict (95\% CI of the paired difference above 0,
Holm-corrected). \textbf{Tier 2 (not reproduced):} a descriptive comparison with our run of the official code under
identical conditions, with no claim against the published number.
\item \textbf{Our four registered variants:} (a) the constraint on the competitor's own network; (b) the constraint on
frozen DINOv2 features of masked images (our pipeline); (c) U-MtE (no artifact labels); (d) the constraint on a small
EfficientNet-B0 (5.3\,M parameters) trained on masked images.
\item \textbf{Our environment (``phase 5''):} ISNet and RR-ClArC trained by us on the thyroid traps with their published
recipes, one selection rule for everyone, registered before any result.
\end{itemize}

\subsection{Results at a glance}
\begin{longtable}{>{\raggedright\arraybackslash}p{4.4cm}>{\raggedright\arraybackslash}p{2.2cm}>{\raggedright\arraybackslash}p{1.3cm}>{\raggedright\arraybackslash}p{6.4cm}}\toprule
Paper (artifact) & Reproduced? & Tier & Outcome\\\midrule\endhead
\textbf{Phase 5}: ISNet, RR-ClArC on thyroid calipers inside the nodule & --- & confirm. & \ok: our constraint beats both, 4/4 comparisons, also better on clean images\\
\textbf{EPVT} (MICCAI 2023 / IEEE TMI 2025; hair, rulers, gel) & yes & 1 & \ok: beaten at both bias levels ($+0.086$, $+0.031$); 0.02--0.04 lower on artifact-free images\\
\textbf{DISC} (ICML 2023; hair) & no & 2 & \ok\ (descriptive): 0.736 vs 0.489 for their code, above on all five trap sets\\
\textbf{RR-ClArC} (Machine Learning 2025; vignette) & no & 2 & \ok\ (descriptive): 0.785 vs 0.524 for their code (0.76 published)\\
\textbf{ISNet} (Nat.\ Commun.\ 2024; class marks) & no & 2 & \ok\ (shows the problem): ISNet ignores corner marks but relies only on marks inside the lungs\\
\textbf{Viviano et al.} (ICLR 2021; tube, location test, 3 seeds) & yes & formal & \mixed: masking's crossover supported; our gain over their model not significant\\
\textbf{Saab et al.} (MLHC 2022; trap sets) & no & 2 & \mixed: tie ($+0.023$, CI $-0.029$ to $+0.076$)\\
\textbf{RR-ClArC} (AAAI 2024; invisible trigger) & yes & 1 & \no: not beaten (the artifact has no location)\\
\textbf{Reveal to Revise} (MICCAI 2023) & no & --- & not compared (released weights did not match the paper)\\
\bottomrule\end{longtable}

\subsection{Our environment: real calipers inside the nodule}
\plain{This is the decisive test: real clinical artifacts, a test set nobody had looked at, and the published methods
run by us under the same rules as ours.}
\begin{table}[H]\centering\small
\begin{tabular}{lccc}\toprule
Method & Trap A (inside) & Trap B (outside) & Clean\\\midrule
Constraint, frozen DINOv2 (ours) & \textbf{0.716} & 0.789 & 0.797\\
Constraint, EfficientNet-B0 (ours) & 0.690 & 0.746 & 0.733\\
ISNet & 0.561 & 0.673 & 0.633\\
RR-ClArC & 0.467 & 0.595 & 0.643\\
U-MtE (ours, no labels) & 0.504 & 0.782 & 0.727\\
ERM (plain training) & 0.440 & 0.546 & 0.625\\\bottomrule
\end{tabular}
\caption{Reversed-test AUROC (three folds). All four registered comparisons are ``beaten'' (e.g.\ DINOv2 vs ISNet
$+0.155$ [$+0.019, +0.316$], vs RR-ClArC $+0.249$ [$+0.192, +0.304$]). ISNet is better with the caliper outside the
nodule than inside it, exactly the location effect.}
\end{table}

\subsection{Their environment}
\begin{itemize}[leftmargin=*]
\item \textbf{EPVT} reproduced after a 37-item audit of its code (one port error fixed; the paper's mixup followed where
the code differs). Our DINOv2 variant beat it at bias 1 (0.782 vs 0.696) and bias 0.9 (0.806 vs 0.775); adding our
constraint to EPVT's own network also improved it. On artifact-free images our frozen variant is 0.02--0.04 lower.
\item \textbf{DISC} did not reproduce. Its released code shares one image cache across trap sets, so later trap sets read
the wrong images; our misaligned runs matched its published baseline (0.368 vs 0.364). Run correctly, DISC scores 0.489;
all our variants are higher (DINOv2 0.736), and also above the published DISC (0.551) and CBDebug (0.580, ICLR 2026).
\item \textbf{ISNet}: with its marks moved from the corner into the lungs, it went from ignoring them (0.84 / 0.83 / 0.83
on own-class / no / swapped marks) to relying on them alone (0.998 / 0.167 / 0.000).
\item \textbf{Viviano et al.}: with a tube inside the lungs+heart ROI, their masked model dropped to 0.615 (0.896 with
the tube outside); our variants scored 0.867--0.876 at every seed, but the registered gain over their standard model was
not significant because at one seed their checkpoint rule kept a model that had not learned the tube.
\end{itemize}

\subsection{Limits found}
\begin{itemize}[leftmargin=*]
\item Class-specific marks inside the ROI (ISNet's design): our variants resist better than ISNet (0.20--0.36 vs 0.03
under swapped marks) but still mostly follow the marks.
\item An invisible trigger spread over the whole image (RR-ClArC, AAAI 2024): not beaten --- outside the method's scope.
\item Segmentation-supervised ERM (Saab et al.) ties with us while using lesion masks as training targets.
\item A trained encoder absorbs strong shortcuts; this is why the frozen encoder is the default.
\end{itemize}

\subsection{Reproducibility findings}
\plain{Five of the eight published results did not reproduce from the authors' code. For three we found concrete causes,
and we emailed the authors of every paper that did not reproduce.}
DISC: shared image cache across trap sets; EPVT: mixup labels paired differently in the code than in the paper; Saab et
al.: the evaluation notebook scores an epoch-1 checkpoint; ISNet: the released model fails even on its own training
sources under the stated preprocessing; RR-ClArC (2025): 0.52 vs 0.76 at the paper's stated settings.
